
% ════════════════════════════════════════════
% 具体分析（5.1.1.具体分析.tex，三级子块）
% ════════════════════════════════════════════
\subsubsection{具体分析}

针对问题一，本问属于确定性线性优化类问题，采用线性规划方法对微网一天的购电与储能充放电策略进行联合优化求解。

本文先对附件一提供的单日分时电价、小区负载与光伏预测数据进行读取与时间对齐，将其整理为一天 144 个 10 分钟时间间隔上的功率序列；依据储能设备的容量、充放电功率与效率约束，将一天内的能量平衡关系与储能荷电状态递推关系写成线性约束；以全天购电费用最小为目标函数，加入储能设备在 0:00 与 24:00 储电量相等的日循环约束，并按附录1给定值固定 0:00 储电量为 6000 kWh，构建标准线性规划模型；最后调用线性规划求解器求得每个时间间隔的购电功率与储能充放电功率，进而得到全天购电量、购电费以及各时段储能充放电量。具体分析流程如图\ref{fig:analysis1_flow}所示。

\begin{figure}[H]
    \begin{center}
        \begin{tikzpicture}[x=1cm, y=0.75cm,
            flowhead/.style={rectangle, rounded corners=3pt, draw=blue!70!black, fill=blue!55!black, text=white, font=\bfseries, minimum width=7.2cm, minimum height=1.1cm, align=center, inner sep=2pt, line width=0.7pt},
            flowbox/.style={rectangle, rounded corners=2.5pt, draw=blue!65!black, fill=blue!12, minimum width=7.2cm, minimum height=1.0cm, align=center, inner sep=2pt, line width=0.6pt},
            flowend/.style={rectangle, rounded corners=2.5pt, draw=blue!75!black, fill=blue!45!black, text=white, font=\bfseries, minimum width=7.2cm, minimum height=1.0cm, align=center, inner sep=2pt, line width=0.7pt},
            flowarrow/.style={thick, -{Stealth}, blue!60!black}
        ]
            \node[flowhead] at (0, 0) (h) {问题一：单日线性规划建模与求解};
            \node[flowbox] at (0, -1.8) (s1) {\textbf{读取单日数据}\\\footnotesize 负荷 · 光伏 · 分时电价};
            \node[flowbox] at (0, -3.6) (s2) {\textbf{时间对齐}\\\footnotesize 整理为 144 个 10 分钟间隔};
            \node[flowbox] at (0, -5.4) (s3) {\textbf{设定决策变量}\\\footnotesize 购电功率与储能充放电功率};
            \node[flowbox] at (0, -7.2) (s4) {\textbf{整理运行约束}\\\footnotesize 储能容量 / 充放电功率 / 效率};
            \node[flowbox] at (0, -9.0) (s5) {\textbf{建立目标函数}\\\footnotesize 全天购电费用最小};
            \node[flowbox] at (0, -10.8) (s6) {\textbf{固定储电初值并施加日循环}\\\footnotesize $E_0=6000$ kWh · SOC 在 0:00 与 24:00 相等};
            \node[flowbox] at (0, -12.6) (s7) {\textbf{构建并求解线性规划}\\\footnotesize 调用线性规划求解器};
            \node[flowend] at (0, -14.4) (s8) {输出最优策略\\\footnotesize 购电量 · 购电费 · 充放电量};
            \draw[flowarrow] (h) -- (s1);
            \draw[flowarrow] (s1) -- (s2);
            \draw[flowarrow] (s2) -- (s3);
            \draw[flowarrow] (s3) -- (s4);
            \draw[flowarrow] (s4) -- (s5);
            \draw[flowarrow] (s5) -- (s6);
            \draw[flowarrow] (s6) -- (s7);
            \draw[flowarrow] (s7) -- (s8);
        \end{tikzpicture}
        \caption{问题一分析流程图}
        \label{fig:analysis1_flow}
    \end{center}
\end{figure}
